\section{本章小结}
\label{sec:rvt-summary}

本章回答的问题是：输入的规律已知、输出由确定的公式算出时，输出服从什么分布。所有计算都从原像规则~\eqref{eq:rvt-preimage-rule}出发：把关于输出的事件写回输入空间，用输入的质量表或联合密度计算。离散时这就是把被送到同一点的质量相加；连续时先得到分布函数，需要密度时再把积分换到输出的坐标上。

换坐标时，记录伸缩的因子各不相同，但道理一样：概率不变，所占的长度或体积变了，密度就按比例反向调整。下表把本章的公式放在一起，最后一列是各自新增的检查。

\begin{BookInlineTable}
  \centering
  \small
  \begin{tabular}{@{}lll@{}}
    \toprule
    映射 & 密度公式 & 须检查的条件 \\
    \midrule
    一元单调 & $f_X(h(y))\lvert h'(y)\rvert$ & 不等号方向、$g'\neq0$、支撑 \\
    一元多分支 & 各逆像贡献之和 & 每支单调、漏支、原子 \\
    同维可逆 & $f_{\bm{X}}(\bm{h}(\bm{y}))\lvert\det D\bm{h}(\bm{y})\rvert$ & 一一、行列式非零、取值区域 \\
    降维 & 补坐标后积掉 & 补出的变换可逆 \\
    和 & $\int f_{X,Y}(x,z-x)\,\mathrm{d}x$ & 独立时才成为卷积 \\
    乘积 & $\int f_{X,Y}(x,z/x)\lvert x\rvert^{-1}\,\mathrm{d}x$ & 零处的原子另计 \\
    比值 & $\int f_{X,Y}(rw,w)\lvert w\rvert\,\mathrm{d}w$ & 分母不取零，尾部与矩 \\
    排序 & $n!\prod_if(x_i)$ & 同值概率为零，否则用计数式 \\
    \bottomrule
  \end{tabular}
  \caption{本章的变换公式与各自的条件。}
  \label{tab:rvt-summary}
\end{BookInlineTable}

几条共同的教训贯穿各节。第一，分布与少数数值特征不同：期望与方差由边缘和协方差决定，完整的输出分布却需要联合规律。第~\ref{sec:rvt-output-distribution}节的三张概率表、第~\ref{sec:rvt-difference-distribution}节的三种依赖都说明了这一点。第二，独立是卷积、封闭性和乘积分解的前提，不能换成“不相关”或“边缘同族”；独立的分组经变换后仍然独立，读同一组输入的两个量则通常不独立。

第三，退化情形要单独处理。导数在一段区间上为零会产生原子，奇异的线性映射使支撑降维，输入在零处的原子会传给乘积，有原子的样本会出现同值。第四，分母靠近零是比值重尾的来源：均匀变量之比与Cauchy分布的均值都不存在，t分布的$k$阶矩存在当且仅当$k<\nu$。

本章顺带补全了前几章留下的几处推导：多元高斯的满维密度来自线性换元，$\chi^2_1$来自平方变换，Rayleigh分布与均匀方向角来自极坐标，Gamma归一化同时给出Dirichlet密度与总量的独立性，t分布来自高斯与$\chi^2$变量之比。

这些计算都把输出的分布写成“输入分布在映射下的像”。第~\ref{chap:probability-spaces}章将在一般的概率空间上统一这种“像分布”：原像规则成为它的定义，本章各种公式成为它在具体映射下的计算，离散、连续与混合的情形也不再分开处理。另一方面，本章的许多变换可以反过来用：$-\ln U$把均匀变量变成指数变量，极坐标把均匀角度和Rayleigh半径变成一对独立高斯，第~\ref{chap:sample-generation}章将利用这些变换从均匀随机数生成指定分布的样本。
